\chapter{LITERATURE REVIEW}
	
This chapter discusses the published academic papers used to support our project. There are two main parts: The first part discussed the technology stack, website development framework, and website development planning that has been implemented and met with success. Such technology stack and planning would then be used as a guideline to our project development. The second part explores any methods or algorithms that could be adapted and used for the project’s main functions development, which would likely be implemented in each function to help fasten the project development.

\section{Design web service academic information system based multiplatform \cite{paper_10}}

This paper proposed a system design that uses Web Services to make academic information available on both web and mobile platforms.

The core of the design is a REST API that uses JSON to exchange data between the server and different types of devices. This approach was chosen because it allows a single database to support many different user interfaces while keeping all the information synchronized in real-time.

The study found that using this web service architecture makes the system much easier to maintain and scale up because the "business logic" on the server is kept separate from how the app looks on the screen.

For future work, the authors suggested adding stronger security protocols to the API and making the database faster so it can handle more students using the app at the same time.

\section{Design and Implementation of a Digital Notice Board and Forum Discussion System \cite{paper_11}}

\section{Course Selling Website Using MERN Stack \cite{paper_12}}

\section{CourseViewer a prototype for visualizing undergraduate courses and student transcripts \cite {paper_15}}

This paper introduces CourseViewer, a prototype that transforms traditional, text-heavy transcripts into interactive visual maps to help students plan their academic paths. This approach is highly relevant for an academic planner because it replaces manual checking with a dynamic system that shows how different courses are connected.

The most useful feature for a modern planner is the graph-based visualization. By representing courses as "nodes" and prerequisites as directed arrows, the system makes it easy to see which classes "unlock" future subjects. This technical design helps prevent enrollment errors by clearly showing if a student has met the necessary requirements before they try to add a class.

Technically, the tool uses the Java language and the Prefuse visualization toolkit to automatically handle the map layout so the user doesn't have to organize it manually. It also uses color-coding—such as green for completed and red for required—which provides an immediate visual status of a student's progress toward graduation.

The study found that these visual maps allow students to identify their remaining credits much faster and more accurately than reading a paper transcript. For future work, the authors suggested adding a "what-if" feature to simulate different majors and integrating the tool directly into online registration systems.


\section{The curriculum prerequisite network: a tool for visualizing and analyzing academic curricula \cite{paper_13}}

This paper treats a university curriculum as a "complex system" that can be mapped using network science to reveal the hidden connections between courses. This is directly applicable to the project’s goal of automating 4-year academic planning for major and general education courses, which currently must be done manually.

The study uses a Directed Acyclic Graph (DAG) where courses are "nodes" and prerequisites are "links". A key technical detail is the identification of "hubs" (courses that unlock many others) and "bridges" (links between different course groups). Implementing this logic allows the interactive planner to visually flag "bottleneck" courses, helping students avoid choices that could delay their graduation by a year.

Technically, the research utilized NetworkX and Gephi to analyze network density and "betweenness centrality". In the project, these metrics can be used to calculate the "long-term impact" of taking a specific class, providing an automated "roadmap" that ensures all prerequisite constraints are met at a glance.


\section{A hyper-heuristic algorithm based on genetic and greedy strategy for university course scheduling problem \cite{paper_3}}
This study addresses the University Course Scheduling Problem (UCSP) by proposing a novel model that balances resources like classrooms and teachers while satisfying complex constraints.

A hyper-heuristic algorithm (GA\_RG\_HH) was developed, combining Genetic Algorithms (GA) with greedy strategies to optimize scheduling efficiency.

Performance was evaluated using real-world datasets, measuring the reduction of soft constraint violations and the efficiency of resource allocation.

The experiment indicated that the proposed GA\_RG\_HH performs better than traditional algorithms like Tabu Search and Particle Swarm Optimization. This is attributed to the algorithm's ability to balance exploration and exploitation through a pre-check strategy and dynamic operator sequences.

Future work includes further optimization of the algorithm for even larger-scale institutional data and the inclusion of more diverse constraint parameters.

\section{A Review of Strategies for Designing, Administering, and Using Student Ratings of Instruction \cite{paper_14}}

